\documentclass[../../main.tex]{subfiles}

\begin{document}

\subsection*{\texttt{genie\_reader}}
  \texttt{genie\_reader} è un dato Complex, Unique e Context, il cui compito
  è rappresentare le strutture dati contenute nel output di GENIE in formato
  rootracker.

  In quanto Complex è responsabile del proprio input; passando un file .root
  nel json di configurazione la sua \texttt{factory} lo apre e cerca un
  \texttt{TTree} di nome ``gRooTracker'' nella cartella ``DetSimPassThru''
  come previsto dalla specifica riportata in figura~\ref{fig:input-structure},
  se il tree è presente viene controllata la versione del formato
  rootracker~\cite[pag. 116]{Andreopoulos_2015} e in base a essa imposta per
  la lettura i soli campi presenti.

  Essendo \texttt{genie\_reader} marcato come Context deve rappresentare i
  dati separati per spill, il formato rootracker però è solo una lista di
  interazioni di neutrini senza alcun concetto di spill.
  Per ovviare a questo problema la factory cerca nel file .root anche il
  \texttt{TTree} ``EDepSimEvents'' che contiene l'output di EDep-Sim diviso
  in spill, e crea una mappa: indice dello spill - range di indici nel
  gRooTracker.


\end{document}